\scanpage{241}
\mspage{121}
% l. 1
that \unsure{$\xi_0$} correspond \inserted{$1-1$} to the
$\xi'\in H^1(S,\mathcal H'_S)$.
% A hatched loop with a descender stands under the H of H'_S (an overwritten letter, perhaps a G).

% l. 2
For such a \deleted{\unsure{$\mathcal H'_S$}} $\xi'$, let
$\delta'(\xi')\in H^2(S,\mu)$
% l. 3
be the obstruction to lifting it to $H^1(S,\unsure{\mathcal G'_S|\mathcal G_S})$,
% l. 4
and $\delta(\xi')=\delta(\xi)$ the obstruction to lifting it
% l. 5
to $H^1(S,\mathcal G_S)$. Unless I am mistaken one has
% l. 6
\[
\delta'(\xi')=\delta(\xi')-\rho_0
\]
% l. 7
i.e.
\[
\delta(\xi')=\delta(\xi)=\delta'(\xi')+\rho_0
\]
% l. 8
and one wants $\delta(\xi)=\rho$ i.e.
% l. 9
\[
\deleted{\unsure{\delta(\xi')}=}\quad
\delta'(\xi')=\rho-\rho_0
\]
% l. 10
and the question again comes down to \deleted{\illegible{3}} \inserted{the}
% l. 11
surjectivity of $H^1(S,\mathcal H'_S)\longrightarrow H^2(S,\mu)$
\deleted{\illegible{1}}
% l. 12
--- I am not getting anywhere. One ought to
% l. 13
consult competent people, like
% l. 14
Giraud or Larry Breen. One feels that one
% l. 15
ought to work with a Gr-stack $\underline{\underline{\mathcal H}}$
% l. 16
of coefficients, with invariants $\underline{\pi}_0\simeq\Gamma$ and
$\underline{\pi}_1\simeq\mu$,
% l. 17
(but not necessarily a Picard stack!)
% l. 18
\deleted{with} and
$H^1(S,\underline{\underline{\mathcal H}})\simeq\text{Classes of maps}$
% l. 19
from $S$ to $B\underline{\underline{\mathcal H}}$, which is a connected
% l. 20
space with $\pi_1(B\underline{\underline{\mathcal H}})\simeq\Gamma$,
$\pi_2(B\underline{\underline{\mathcal H}})=\mu$.
% l. 21
Here the \unsure{Postnikoff} class in $H^3(\Gamma,\mu)$ is zero,
\margin{240}

\scanpage{242}
% l. 1
[but perhaps there is no reason to
% l. 2
cling too much to this hypothesis,
% l. 3
corresponding to the existence of an \inserted{extension} $\mathcal G$ of $\mathcal H$
\inserted{by $\pi_1(\mathcal H^\circ)$} \deleted{\illegible{1}}
% l. 4
which gives back the universal extension of
% l. 5
$\mathcal H^\circ$ by $\pi_1(\mathcal H^\circ)$] one has a \unsure{homomorphism}
% l. 6
$B\mathcal H$ \deleted{\illegible{1}} $\longrightarrow B\underline{\underline{\mathcal H}}$
which induces an isomorphism for
% The struck mark stands under the first BH (perhaps a cancelled double underline).
% l. 7
the $\pi_i$ for $i\leq2$, and for $i\leq n$
(where $n\geq2$ \inserted{is given})
% l. 8
iff $\pi_i(B\mathcal H)=0$ for $2<i\leq n$, i.e.
% l. 9
$\pi_i(\mathcal H)=0$ for $2\leq i\leq n-1$ [in the case
% l. 10
we are interested in $\mathcal H=A_0$, one can
% l. 11
take $n=3$). This implies that for
% l. 12
every C-W-complex $S$ of dimension $\leq n$,
% l. 13
one has \inserted{(Hot denoting the morphisms in
the unpointed homotopy category)}
% l. 14
\[
\mathrm{Hot}(S,B\mathcal H)\xrightarrow{\sim}
\mathrm{Hot}(S,B\underline{\underline{\mathcal H}})
\]
% l. 15
\unsure{i.e.|i.e.}
\[
H^1(S,\mathcal H)\xrightarrow{\simeq}
H^1(S,\deleted{\unsure{B}}\underline{\underline{\mathcal H}})
\]

% l. 16
But if the \unsure{Postnikoff}-Sinh invariant
% l. 17
$\unsure{K|k}\in H^3(\Gamma,\mu)$ is zero,
\deleted{\illegible{1}} (as well as the action
% l. 18
of $\Gamma$ on $\mu$) then unless I am mistaken
$B\underline{\underline{\mathcal H}}$ is iden%
% l. 19
tified with a product $K(\Gamma,1)\times K(\pi,2)$,
\deleted{and}
% l. 20
(this ``identification'' depending precisely on the
\margin{241}

\scanpage{243}
\mspage{122}
% l. 1
choice of the $H^2(\pi_1,\pi_2)\simeq H^2(\Gamma,\mu)$ !), and
% l. 2
hence
\[
H^1(S,\underline{\underline{\mathcal H}})\xrightarrow{\simeq}
H^1(S,\Gamma)\times H^2(S,\mu)
\]
% l. 3
for every space $S$, hence for $S$ a C-W-%
% l. 4
complex of dimension $\leq n$
% l. 5
\[
(71)\qquad
H^1(S,\mathcal H)\xrightarrow{\simeq}H^1(S,\Gamma)\times H^2(S,\mu)
\]

% l. 6
So the classification of bundles of topological
% l. 7
spheres over a topological space $S$,
% l. 8
for a C-W complex $S$ of dimension $\leq3$,
% l. 9
really does work. Of course, the
% l. 10
fact that one is obliged here to restrict oneself to
% l. 11
\deleted{\illegible{1}} $S$ of dimension $\leq3$ is due to the fact
% l. 12
that we have found \deleted{\unsure{only}}
(via $\underline{\underline{\mathcal H}}$)
% l. 13
only a $2$-isotopic (and not
% l. 14
$\infty$-isotopic) description of the category of
% l. 15
topological $2$-spheres.

% l. 16
I would like to take up the classification again
% l. 17
for an arbitrary CW-complex $S$,
% l. 18
using the fact that the inclusion
% l. 19
\[
\widetilde{\mathrm{GP}}(1,\mathbb C)\hookrightarrow A_0
\]
\margin{242}

\scanpage{244}
% l. 1
is a homotopy equivalence, and therefore
% l. 2
induces a bijection
% l. 3
\[
H^1(S,\widetilde{\mathrm{GP}}(1,\mathbb C))\longrightarrow H^1(S,A_0)
\]
% l. 4
and even an $\infty$-equivalence of the $\infty$-categories of
% l. 5
\deleted{bundles} \inserted{torsors} over $S$, with \deleted{\illegible{1}} groups
$\widetilde{\mathrm{GP}}$ (corresponding
% l. 6
to bundles of conformal spheres, or of
% l. 7
projective lines over an \deleted{field} $\mathbb R$-algebra
\inserted{(unspecified)}
% l. 8
isomorphic to $\mathbb C$) and with group $A_0$ --- corresponding
% l. 9
to sphere bundles over $S$.

% l. 10
The first invariant of a bundle with group $\widetilde{\mathrm{GP}}$
% l. 11
is a \deleted{bundle} torsor \unsure{$\eta$|$T$} under $(\mathbb Z/2\mathbb Z)_S$
\deleted{\illegible{2}} defined (\unsure{up to isomorphism}
% l. 12
\uline{not unique}\inserted{) by a} $\chi\ (\in H^1(S,\mathbb Z/2\mathbb Z))$.
% Physical order: "non unique" opens l. 12 under "est un"; "près) par un" is raised after the
% parenthesis of H^1; linearised as the sense requires.
\margin{i.e. an \unsure{étale covering|double covering} of degree $2$.
The datum of \unsure{$\eta$|$\chi$} amounts to the datum of a system \unsure{$T$} of twisted integers
over $S$.}
% l. 13
\inserted{This torsor} will serve to twist $\mathbb C$
(via the \inserted{faithful} operation of $\mathbb Z/2\mathbb Z$ on
% l. 14
the $\mathbb R$-extension $\mathbb C$), \deleted{and} whence a locally
\deleted{\illegible{1}}
% l. 15
constant bundle of \deleted{fields} \inserted{quadratic} extensions
of $\mathbb R$ \deleted{\illegible{1}},
% l. 16
say $C$, and one \deleted{will take the} will enrich $S$
% l. 17
with the sheaf \inserted{$\underline C_S$ (or $\underline C$)} of continuous sections of
% l. 18
$C$ (which is a way of twisting
% l. 19
\deleted{$\mathrm{Sl}(2,$} $\mathrm{GP}(1,\mathbb C)_S$
into $\mathrm{GP}(1,C)_S=\mathrm{GP}(1,\underline C)$))
% l. 20
and the point is to \deleted{\unsure{therefore}} describe in an
% l. 21
understandable way \deleted{the category} \inserted{--- passing} if need be,
% l. 22
to the \deleted{\unsure{$\infty$}}$n$-isotopic ones, for $n=\ldots$ --- the category
\margin{243}

\scanpage{245}
\mspage{123}
% l. 1
--- depending on the choice of $\chi$ via $\underline C$ --- of
% l. 2
Azumaya algebras of rank $4$ over $\underline C$,
% l. 3
or again of bundles of projective lines
% l. 4
over $\underline C$, or of torsors under $\mathrm{GP}(1,\underline C)$.

% l. 5
Now using \deleted{$\mathrm{Sl}(2,\underline C)\longrightarrow
\mathrm{GP}(1,\underline C$} \inserted{the \unsure{exact} sequence}
% l. 6
\[
1\longrightarrow\mu_{2S}\longrightarrow\mathrm{Sl}(2,\underline C_S)
\longrightarrow\mathrm{GP}(1,\underline C_S)\longrightarrow1
\]
% l. 7
one associates to an object of the category
% l. 8
--- say a torsor under $\mathrm{GP}(1,\underline C_S)$ --- the
% l. 9
\inserted{fibered (over variable open subsets of $S$)} category of these liftings to
$\mathrm{Sl}(2,\underline C_S)$,
% l. 10
which is a \uline{gerbe} \inserted{$\mathcal G$} banded by $\mu_2$.
% l. 11
It is this gerbe that is the second
% l. 12
complete invariant --- finer than
% l. 13
its class \deleted{in $H^2($} \inserted{up to equivalence, which is identified with a}
% l. 14
$\xi\in H^2(S,\mu_2)=H^2(S,\{\pm1\})$
(playing the role
% l. 15
of a Brauer group)

% l. 16
An isomorphism of bundles, with invariants
% l. 17
$(T,\mathcal G)$, $(T',\mathcal G')$, will define an isomorphism $T\simeq T'$
% l. 18
(whence an isomorphism \deleted{$\mathrm{GP}(1,C$}
$\underline C_T\simeq\underline C_{T'}$ whence
% l. 19
an isomorphism $\mathrm{GP}(1,\underline C_T)\simeq\mathrm{GP}(1,\underline C_{T'})$) --- and
% l. 20
an equivalence of gerbes --- one should
\margin{244}

\scanpage{246}
% l. 1
spell out that for two isomorphisms
% l. 2
$f,g$ of the torsors, \deleted{and} whence
% l. 3
\diagram{\begin{tikzcd}[ampersand replacement=\&,column sep=3.4em]
f_T,g_T:\ T
 \arrow[r,shift left=2.6ex,line width=0.7pt,"\unsure{\simeq}"]
 \arrow[r,phantom,"\unsure{\deleted{\xrightarrow{\sim}}}" description]
 \arrow[r,shift right=2.6ex,"\sim"]
\& T'
\& f_{\mathcal G},g_{\mathcal G}:\mathcal G
 \arrow[r,shift left=1.6ex,"\unsure{\simeq}"]
 \arrow[r,shift right=1.6ex,"\unsure{\simeq}"]
\& \mathcal G'
\end{tikzcd}}
% l. 4
and for every homotopy $h_t$ $(0\leq t\leq1)$ from
% l. 5
$f$ to $g$, one finds $f_T=g_T$ and an
% l. 6
isomorphism $h_*:f_{\mathcal G}\xrightarrow{\sim}g_{\mathcal G}$ of equivalences of
% l. 7
gerbes, which depends only on the homotopy
% l. 8
class of this homotopy. One
% l. 9
thus finds a $2$-functor, from the
% l. 10
$2$-category of \deleted{bundles} torsors over $S$
% l. 11
with group $\widetilde{\mathrm{GP}}(1,\mathbb C)$, to the $2$-category
% l. 12
formed by pairs \inserted{$(T,\mathcal G)$} of a
\deleted{torsor} \inserted{system} of %
% l. 13
twisted integers, \deleted{and} ($\Longleftrightarrow$ of a torsor over
% l. 14
$S$ with group \deleted{\illegible{1}} \inserted{$\mathbb Z/2$}) and of a
% l. 15
\deleted{\unsure{$\{\pm1\}=$}} $\mu_2(\mathbb C)$-gerbe $\mathcal G$ over $S$. This
% l. 16
$2$-functor, unless I am mistaken, is $2$-faithful
% l. 17
under the conditions spelled out
% l. 18
above (\deleted{\illegible{2}} (dimension $S\leq4$), and
% l. 19
is $3$-faithful (i.e. the injective map
% l. 20
$H^1(S,\widetilde{\mathrm{GP}})\unsure{\longrightarrow|{-}}
H^1(S,\mathbb Z/2\mathbb Z)\times H^2(S,\mu_2)$ is surjective)
\margin{245}

\scanpage{247}
\mspage{124}
% l. 1
if one even has dimension $S\leq3$. One expects it
% l. 2
to be $1$-faithful as soon as dimension $S\leq5$, $0$-faithful
% l. 3
as soon as dimension $S\leq6$.
\margin{??}
% A vertical marginal stroke accompanies approximately lines 1--3.
[\uline{N.B.} \inserted{the assertion} ``$0$-faithful'' means
% l. 4
that if \inserted{above} one has two homotopies $h,h'$ from $f$ to $g$,
% l. 5
such that $h_*=h'_*,$ then $h$ and $h'$ are homotopic.
% l. 6
The condition ``$1$-faithful'' means \inserted{$0$-faithful and} that every
% l. 7
isomorphism from $f_{\mathcal G}$ to $g_{\mathcal G}$
\deleted{\unsure{comes from a}} is of the
% l. 8
form $h_*$ (with $h$ well determined up to homotopy,
% l. 9
by the preceding condition of $0$-faithfulness). The
% l. 10
$2$-faithful condition means that \inserted{in addition, for} every
\deleted{equivalence}
% l. 11
\deleted{\unsure{given}} $\varphi:T\simeq T'$,
$\psi_{\mathcal G}:\mathcal G\xrightarrow{\approx}\mathcal G'$,
there exists
% l. 12
an isomorphism of bundles $f$ such that $\varphi=f_T$,
\deleted{$\psi$ \illegible{1}}
% l. 13
and an isomorphism $\psi\simeq f_{\mathcal G}$. Finally $3$-faithful means
% l. 14
that in addition, for every pair $T,\mathcal G$,
% l. 15
$\exists\ T',\mathcal G'$ coming from a bundle, \deleted{and of}
% l. 16
and an isomorphism $T\simeq T'$ and an equivalence
% l. 17
$\mathcal G\approxeq\mathcal G'$. These categorical notions
% l. 18
should simply interpret that, if
% l. 19
one looks at $\mathcal H=\widetilde{\mathrm{GP}}(1,\mathbb C)
\xrightarrow{\alpha}\underline{\underline{\mathcal H}}$
--- where $\underline{\underline{\mathcal H}}$
% l. 20
is deduced from $\mathcal H$ by ``killing the homotopy \unsure{gr\textsuperscript{o}ups}
% l. 21
$\pi_i$ for $i\geq2$'', so that
\margin{246}

\scanpage{248}
% l. 1
$\alpha$ induces an isomorphism of the $\pi_i$ for $i\leq1$
% l. 2
(and even for $i\leq n$ if
$\pi_i(\mathcal H)=0$ for $1<i\leq n$)
% l. 3
\deleted{and $B_\alpha$ an} whence
% l. 4
\[
B_\alpha:B\mathcal H\longrightarrow B\underline{\underline{\mathcal H}}
\]
\margin{\uline{N.B.} $B\underline{\underline{\mathcal H}}$ is deduced from
$B\mathcal H$ by killing the $\pi_i$ for $i\geq3$.}
% l. 5
(inducing an isomorphism of the $\pi_i$ for $i\leq n+1$),
% l. 6
then taking for \deleted{every} \inserted{a given space} $S$ the
% l. 7
corresponding map of spaces of continuous maps
% l. 8
\[
\deleted{\mathrm{Hot}}
\underline{\mathrm{Cont}}(S,B\mathcal H)
\unsure{\longrightarrow|{-}}
\underline{\mathrm{Cont}}(S,B\underline{\underline{\mathcal H}})
\]
% l. 9
this induces an isomorphism for the
% l. 10
$\pi_i$ $(0\leq i\leq2)$. \deleted{\unsure{Since}} It would seem that
% l. 11
$0$-faithfulness is a condition of injectivity
% l. 12
for the $\pi_2$, $1$-faithful bijectivity for the
% l. 13
$\pi_2$ and injectivity for the $\pi_1$, $2$-faithful
% l. 14
bijectivity for $\pi_2$ and $\pi_1$ and injectivity
% l. 15
for $\pi_0$, $3$-faithful bijectivity for $\pi_0,\pi_1,\pi_2$.

% l. 16
Since $\pi_i(\mathrm{Cont}(S,B))\simeq
\pi_0(S\times S^i,B)$,
% l. 17
\deleted{one sees} \inserted{it would seem}
that one is led, for $0$-%
% l. 18
faithfulness, to make the draco%
% l. 19
nian hypothesis that $S\times S^2$ of dimension $\leq4$, for
% l. 20
$1$-faithfulness that $S\times S^2$ of dimension $\leq3$ i.e. $S$
% l. 21
of dimension $\leq1$ ? (\unsure{This|That} would then imply
% l. 22
$3$-faithfulness\ldots)
\margin{247}

\scanpage{249}
\mspage{125}
% l. 1
Let us take the arithmetic analogue, of a description
% l. 2
(more or less ``faithful'') \deleted{by \illegible{1}}
of a ``profinite'' nature
% l. 3
of the projective lines \deleted{defined} (possibly
% l. 4
twisted) defined over a finitely generated
% l. 5
field $K$ \inserted{over $\mathbb Q$ (more generally over an arbitrary scheme $S$)},
these a priori form a category
% l. 6
and nothing more --- a groupoid (not necessarily
% l. 7
connected) I do not know how to make a
% l. 8
$2$-category out of them reasonably, because for two
% l. 9
isomorphisms $f,g$ \inserted{$(=f\circ u)$} of such torsors, to define the
% l. 10
homotopisms from $f$ to $g$, i.e. for a ``form''
% l. 11
$G$ of $\mathrm{GP}(1)_S$ and a section $u$ of $G$,
% l. 12
to define the ``homotopies'' from $u$ to the iden%
% l. 13
tity, \deleted{and the \inserted{homotopy} classes of such \illegible{1}}
\inserted{or rather} a
% l. 14
notion which replaces the notion of a homo%
% l. 15
topy class of such homotopies. \deleted{But} \inserted{This would be}
% l. 16
probably a lifting of $u$ to a
\deleted{
% l. 17
\unsure{If} $S$ is \illegible{1} of a geometric point
% l. 18
\illegible{1} \unsure{every section} defines a \illegible{1}
% l. 19
\unsure{if $K$ is algebraic} For an algebraically
% l. 20
closed base field \unsure{at least} (the \unsure{sections}
% l. 21
defining a geometric point) one
% l. 22
can think of taking the ``classes of
% l. 23
paths from $f$ to \illegible{1} from $u$ to
}
% l. 24
section of the \deleted{\unsure{$\widetilde{\mathrm{GP}}$}}
simply
% l. 25
connected covering $\widetilde G$ of $G$ !
\margin{248}

\scanpage{250}
% l. 1
I have something rather amusing here,
% l. 2
but which I am not going to pursue --- in
% l. 3
\unsure{any case|any case} it is clear that the ``description''
% l. 4
of the projective lines (over a finitely generated
% l. 5
field, say) that one ends up with,
% l. 6
has nothing faithful about it --- not even $0$-faithful !
% l. 7
Thus all the automorphisms of $\mathbb P^1_K$ coming from
% l. 8
$\mathrm{Sl}(2,K)$ will be identified with the identity ---
% l. 9
that is a bit brutal ! But I
% l. 10
realize that the conceptual work
% l. 11
around the ``Brauer'' themes is not
% l. 12
finished --- that there are things to be understood
% l. 13
for the study of the category
% l. 14
(\inserted{which should be} a \uuline{2}-category !)
of Azumaya Algebras
% l. 15
\deleted{\unsure{of}} \uline{fixed} degree (here $4$)\ldots
\margin{249}